\documentclass{article}
\usepackage{amsmath,amssymb}
\usepackage[margin=1in]{geometry}
\usepackage{xurl}
\usepackage[hidelinks]{hyperref}
\setlength{\emergencystretch}{2em}
\clubpenalty=10000
\widowpenalty=10000
% Shared commands for notes in docs/paper-gaps/.

\DeclareUnicodeCharacter{2080}{\ifmmode{}_{0}\else\textsubscript{0}\fi}
\DeclareUnicodeCharacter{2081}{\ifmmode{}_{1}\else\textsubscript{1}\fi}
\DeclareUnicodeCharacter{2082}{\ifmmode{}_{2}\else\textsubscript{2}\fi}
\DeclareUnicodeCharacter{2099}{\ifmmode{}_{n}\else\textsubscript{n}\fi}

\newcommand{\Lean}{\textsc{Lean}}
\ExplSyntaxOn
\NewDocumentCommand{\leanid}{m}{
  \ifmmode
    \textsf{\detokenize{#1}}
  \else
    \group_begin:
    \urlstyle{sf}
    \tl_set:Nn \l_tmpa_tl {#1}
    \tl_replace_all:Nnn \l_tmpa_tl {\_} {_}
    \exp_args:NV \path \l_tmpa_tl
    \group_end:
  \fi
}
\ExplSyntaxOff
\providecommand{\tr}{\operatorname{tr}}
\newcommand{\tnleanIssueBase}{https://github.com/LionSR/TNLean/issues/}
\newcommand{\tnleanPrBase}{https://github.com/LionSR/TNLean/pull/}
\newcommand{\tnleanDocBase}{https://sirui-lu.com/TNLean/docs/}
\newcommand{\ghissue}[1]{\href{\tnleanIssueBase#1}{GitHub issue \##1}}
\newcommand{\ghpr}[1]{\href{\tnleanPrBase#1}{PR \##1}}
\newcommand{\doclink}[1]{\href{\tnleanDocBase#1.html}{\path{#1}}}

% Dirac notation and the identity operator, as in blueprint/src/macros/common.tex.
% \providecommand keeps note-local definitions authoritative.
\usepackage{dsfont}
\providecommand{\ket}[1]{|#1\rangle}
\providecommand{\bra}[1]{\langle #1|}
\providecommand{\braket}[2]{\langle #1 | #2 \rangle}
\providecommand{\ketbra}[2]{|#1\rangle\!\langle #2|}
\providecommand{\Id}{\mathds{1}}
\providecommand{\one}{\mathds{1}}
\providecommand{\C}{\mathbb{C}}
\providecommand{\MN}[1]{M_{#1}(\C)}
\providecommand{\MD}{M_D(\C)}

% Verdict marker: \gapnote{<kind>}{<status>}. Kind is the policy
% classification (clarification, local-correction, scope-restriction,
% unfaithful, false-source, open-gap); status is open, wip (elimination
% underway), resolved, or historical. Machine-read by the site index;
% prints a verdict line.
\newcommand{\gapnote}[2]{\par\noindent\textbf{Verdict:} #1 (#2).\par}

\title{Ground-Space Subtraction in the Infinite-Volume Gap Argument}
\author{}
\date{2026-10-02}
\begin{document}
\maketitle
\gapnote{open-gap}{wip}

\section{The source conclusion}
Nachtergaele's gap theorem \cite{Nachtergaele1996Spectral} includes an
infinite-volume conclusion. Its hypotheses include a positive finite-range
interaction, equality of its zero-energy state face with the prescribed
GVBS face, and equality of the open-chain kernel with the prescribed
local support space at every sufficiently large length. The interaction
range need not exceed a separately supplied injectivity length.
For each pure GVBS ground state $\eta$ and local observable $X$ with $\eta(X)=0$, the uniform finite-volume gap
$\gamma>0$ also satisfies
\begin{align}
    \lim_{M\to-\infty,\,N\to+\infty}
        \eta\bigl(X^*[H_{[M,N]},X]\bigr)
        &\geq\gamma\eta(X^*X).
    \label{eq:infinite_gap_source}
\end{align}
The local source is \path{References/cond-mat_9410110/main.tex},
lines 933--947. Its final proof, at lines 2649--2675, uses convergence
of finite-volume pure ground-state expectations to those of one pure
infinite-volume sector.

\section{The finite-volume estimate}
Let $H\geq0$ act on a finite-dimensional Hilbert space, let
$K=\ker H$, and suppose $H\geq\gamma P_{K^\perp}$. For every vector
$v$, orthogonal decomposition gives
\begin{align}
    \langle Hv,v\rangle
        &\geq\gamma\bigl(\|v\|^2-\|P_Kv\|^2\bigr).
    \label{eq:infinite_gap_subtraction}
\end{align}
If $\psi\in K$, then
$\langle\psi,X^*[H,X]\psi\rangle=\langle HX\psi,X\psi\rangle$.
Consequently
\begin{align}
    \langle\psi,X^*[H,X]\psi\rangle
        &\geq\gamma\bigl(\|X\psi\|^2-\|P_KX\psi\|^2\bigr).
    \label{eq:infinite_gap_commutator}
\end{align}
These assertions are proved in
\path{TNLean/MPS/ParentHamiltonian/Martingale/GroundExpectationGap.lean}.

The projection in (\ref{eq:infinite_gap_commutator}) is onto the entire
ground space. For a normalized ground vector $\psi$, the condition
$\langle\psi,X\psi\rangle=0$ removes the component along $\psi$;
it does not remove components along other vectors in $K$. Thus scalar
centering alone does not justify replacing the projection term by zero
in a degenerate finite volume.

\section{The proved limiting implication}
Consider Hilbert spaces and operators indexed by expanding volumes $n$.
Suppose their Hamiltonians have the same gap $\gamma$ for all sufficiently
large $n$, and let $\psi_n\in\ker H_n$. If
\begin{align}
    \|X_n\psi_n\|^2&\longrightarrow a,
    \label{eq:infinite_gap_norm_limit}\\
    \|P_{\ker H_n}X_n\psi_n\|^2&\longrightarrow0,
    \label{eq:infinite_gap_projection_limit}\\
    \langle\psi_n,X_n^*[H_n,X_n]\psi_n\rangle&\longrightarrow e,
    \label{eq:infinite_gap_energy_limit}
\end{align}
then (\ref{eq:infinite_gap_commutator}) implies $e\geq\gamma a$.
The formal limiting assertion permits the Hilbert space to vary with $n$.
Its projection-decay hypothesis is explicit; it is not counted as a
formalization of (\ref{eq:infinite_gap_source}).

When $\ker H_n=\mathbb C\psi_n$ and $\|\psi_n\|=1$ for all sufficiently
large $n$, one has
\begin{align}
    \|P_{\ker H_n}X_n\psi_n\|^2
        &=|\langle\psi_n,X_n\psi_n\rangle|^2.
    \label{eq:infinite_gap_rank_one_projection}
\end{align}
Thus convergence of the scalar mean to zero suffices in this case.
The corresponding finite-volume estimate and limiting implication are
also proved in the cited module. Finite-volume uniqueness is an explicit
hypothesis; it is not inferred from purity of the limiting state.

\section{Faithful generating data and finite supports}

Let $A$ be a tensor of positive bond dimension and let $\rho>0$. For a
$k$-site matrix $X$, put
\[
 \eta_{A,\rho}(X)=\frac{1}{\operatorname{tr}\rho}
   \sum_{\sigma,\tau}X_{\sigma\tau}
       \operatorname{tr}(A^\tau\rho(A^\sigma)^*).
\]
No stationarity or normalization is required for the following equivalence:
\[
 \eta_{A,\rho}(X^*X)=0\quad\Longleftrightarrow\quad
 G_A(k)\subseteq\ker X.
\]
Indeed, writing $C_b=\sum_\sigma X_{b\sigma}A^\sigma$, one has
\[
 (\operatorname{tr}\rho)\eta_{A,\rho}(X^*X)
   =\operatorname{tr}\!\left(\rho\sum_b C_b^*C_b\right).
\]
Faithfulness of $\rho$ makes this trace zero precisely when every $C_b$ is
zero. The identity
$(X\Gamma_k(Y))_b=\operatorname{tr}(C_bY)$ and the nondegeneracy of the
trace pairing finish the argument. This is proved in
\path{TNLean/MPS/ParentHamiltonian/LocalFaithfulSupport.lean}.
The argument applies also at $k=0$.

If $\sum_iA_i^*A_i=I$, the functional is positive and normalized.
It therefore has a positive representing matrix $\sigma_k$ with trace one:
\[
 \eta_{A,\rho}(X)=\operatorname{tr}(\sigma_kX),
 \qquad \operatorname{ran}\sigma_k=G_A(k).
\]
The range equality follows by comparing the vanishing conditions for
$\eta_{A,\rho}(X^*X)$ and
$\operatorname{tr}(\sigma_kX^*X)$, then testing the orthogonal-complement
projections of the two subspaces. The result is proved in
\path{TNLean/MPS/ParentHamiltonian/LocalDensitySupport.lean}.
Stationarity is unnecessary for these finite densities. Without stationarity
the conclusion does not assert consistency across intervals.

These finite-volume statements take the faithful generating data as given.
Nachtergaele passes from an arbitrary GVBS presentation to faithful minimal
generating data \cite{Nachtergaele1996Spectral}, local source lines 1724--1738.
That passage is not formalized here; the support results therefore carry the
faithful generators as an explicit hypothesis. Removing it requires
formalizing the minimal-data construction from a general presentation.
Likewise, results stated for a primitive tensor with canonical data, for a
normalized periodic tensor, or for a finite periodic family take that tensor
presentation as given. Deriving such a presentation from a general GVBS
boundary-limit presentation is part of the same unformalized passage, and the
same elimination plan applies.

\section{The primitive-sector quasi-local inequality}
For one normalized primitive tensor with a positive-definite invariant
matrix, the tensor-specific projection decay is now proved. The exact
boundary compression of an interior observable is the reshuffling of
$E_A^\ell E_X E_A^r$. When both free intervals tend to infinity, this
converges to $\eta_A(X)$ times the limiting boundary Gram operator.
These identities and limits are proved in
\path{TNLean/MPS/ParentHamiltonian/LocalObservableInsertion.lean}.
The required general compression inequality is proved in
\path{TNLean/MPS/ParentHamiltonian/GroundSpaceCompressionNorm.lean}:
for an injective boundary map $T$, with $K=(T^*T)^{-1}$ and range
projection $P$, it gives
\begin{align}
    \|PXP-cP\|&\leq\|K\|\,\|T^*XT-cT^*T\|.
    \label{eq:infinite_gap_compression_bound}
\end{align}
An eventual uniform bound on $K$ transfers vanishing of the pullback defect
to vanishing of the full compression defect, even when the Hilbert spaces
vary. In particular, for $\eta_A(X)=0$, the complete ground-space projection
of $X\psi_n$ tends to zero uniformly over unit boundary ground vectors.
This consequence and local expectation convergence are proved in
\path{TNLean/MPS/ParentHamiltonian/BoundaryObservableCompression.lean}.
The finite expectations are positive, normalized, complex linear, and
consistent under adjoining identity sites. Their operator norm is one;
these state properties are proved in
\path{TNLean/MPS/ParentHamiltonian/LocalObservableExpectation.lean} and
\path{TNLean/MPS/ParentHamiltonian/LocalObservableState.lean}.

The fixed local commutator identity is also proved. For an interaction $h$
of range $R\geq1$ and an observable $X$ on $k\geq1$ sites, put
$t=R-1$, $w=2t+k$, and $B=I_t\otimes X\otimes I_t$. If both free
intervals have at least $t$ sites, then $X_n^*[H_n(h),X_n]$ is the
identity-tensor inclusion of $B^*[H_w(h),B]$. This locality identity is in
\path{TNLean/MPS/ParentHamiltonian/BulkObservableCommutator.lean}.
Combining it with the primitive expectation and full projection limits gives
\begin{align}
    \gamma\operatorname{Re}\eta_A(X^*X)
        &\leq\operatorname{Re}\eta_A\bigl(B^*[H_w(h),B]\bigr).
    \label{eq:primitive_local_commutator_gap}
\end{align}
For every positive interaction $h$ with kernel $G_R(A)$, a positive
$\gamma$ independent of $X$ is derived at every $R\geq D^4+1$, where
$D>0$ is the bond dimension. Comparison with the canonical interaction
preserves the open-chain kernel and gives a uniform gap. There is no
supplied energy-limit assumption and no finite-volume uniqueness assumption.
The canonical and positive-interaction statements are proved in
\path{TNLean/MPS/ParentHamiltonian/PrimitiveLocalCommutatorGap.lean} and
\path{TNLean/MPS/ParentHamiltonian/PrimitiveLocalParentInteractionGap.lean}.

The consistent local expectations now have a continuous positive extension
$\omega_A$ to the quasi-local algebra, with $\omega_A(I)=\|\omega_A\|=1$.
Finite-region compatibility, interval evaluation, and translation invariance
are proved in
\path{TNLean/MPS/ParentHamiltonian/LocalObservableRegionExpectation.lean},
\path{TNLean/MPS/ParentHamiltonian/LocalObservableQuasiLocalState.lean}, and
\path{TNLean/MPS/ParentHamiltonian/LocalObservableTranslationInvariance.lean}.
If $j_{a,k}$ includes a $k$-site matrix observable at position $a$, then
$\omega_A(j_{a,k}(X))=\eta_A(X)$. Consequently
(\ref{eq:primitive_local_commutator_gap}) is also the quasi-local inequality
$\gamma\operatorname{Re}\omega_A(j_{a,k}(X)^*j_{a,k}(X))
\leq\operatorname{Re}\omega_A(j_{a-t,w}(B^*[H_w(h),B]))$.
This formulation is proved in
\path{TNLean/MPS/ParentHamiltonian/QuasiLocalPrimitiveGap.lean}.

\newpage
\section{The multiblock projection limit and sector purity}
For several sectors, the ground-space projection also contains the other
sectors. Their local means need not vanish when $\eta_A(X)=0$ in the
chosen sector. These cross-sector components are now removed without
asserting scalarity on the entire multiblock ground space. A ground vector
in sector $\alpha$ lies in its prefix support. A distant observable commutes
with the prefix projections, so
\begin{align}
    \|P_\beta(N)X_{\ell,r}P_\alpha(N)\|
        &\leq\|P_\beta(\ell)P_\alpha(\ell)\|\,\|X\|
          \longrightarrow0 \qquad (\alpha\neq\beta).
    \label{eq:multiblock_prefix_cross_compression}
\end{align}
This estimate is proved in
\path{TNLean/MPS/ParentHamiltonian/BoundaryCrossObservableCompression.lean}.
The diagonal compression tends to zero by centering. The sum of the sector
projections approaches the projection onto their joint span. Thus all
column compressions vanish in norm, and so does the full ground projection
of $X\psi_n$ for bounded vectors in the chosen sector. The finite-sector
assembly and tensor consequence are proved in
\path{TNLean/MPS/ParentHamiltonian/GroundSpaceSectorCompressionDecay.lean} and
\path{TNLean/MPS/ParentHamiltonian/BlockBoundaryObservableCompression.lean}.

For a weighted direct sum of normalized primitive blocks with nonzero
weights, simultaneous injectivity at $S>0$ gives the exact joint kernel
and a uniform gap at every $R\geq S+1$. Every positive interaction with
this local kernel consequently has one positive commutator constant for
all sectors and centered local observables. The finite-interval statement
and its formulation in the constructed quasi-local sector states are proved
in
\path{TNLean/MPS/ParentHamiltonian/MultiblockLocalCommutatorGap.lean} and
\path{TNLean/MPS/ParentHamiltonian/QuasiLocalMultiblockGap.lean}.
No energy limit, projection limit, or finite-volume uniqueness is supplied
as a hypothesis of these existence results. A further corollary starts with
the original simultaneous word-spanning family and constructs primitive
representatives with faithful invariant matrices, preserving every block
ground space and the original positive interaction.

The exact local support of the completed state is also proved in
\path{TNLean/MPS/ParentHamiltonian/LocalParentExpectation.lean}:
an observable annihilating $G_k(A)$ has expectation zero at every position.
In particular, $\omega_A(j_{a,N}(P_{G_N(A)}))=1$ for every interval,
including $N=0$. For any competing normalized positive state $\varphi$
with the same support constraints, Cauchy--Schwarz gives
$\varphi(PXP)=\varphi(X)$ on each interval. Primitive scalar compression
on expanding intervals then identifies all its local expectations with those
of $\omega_A$. Continuity proves $\varphi=\omega_A$ on the completed
algebra. In a convex decomposition of $\omega_A$, positivity forces both
constituents to inherit each support constraint. Uniqueness therefore proves
purity among all states, without imposing translation invariance on the
constituents. These statements are proved in
\path{TNLean/MPS/ParentHamiltonian/PrimitiveQuasiLocalUniqueness.lean} and
\path{TNLean/MPS/ParentHamiltonian/PrimitiveQuasiLocalPurity.lean}.

The literal quasi-local commutator observable is eventually equal to the
fixed patch, as proved in
\path{TNLean/MPS/ParentHamiltonian/QuasiLocalCommutatorLocality.lean}.
Local support and positivity make its sector expectation real and
nonnegative, by
\path{TNLean/MPS/ParentHamiltonian/LocalCommutatorExpectation.lean}.
Consequently the literal limit in (\ref{eq:infinite_gap_source}) now exists
and satisfies the complex-order inequality in each constructed pure sector,
with one common positive constant at $R\geq S+1$. This is proved in
\path{TNLean/MPS/ParentHamiltonian/QuasiLocalMultiblockLimitGap.lean}.
The zero-energy face is now classified for a given finite family of
inequivalent primitive sectors under the eventual kernel hypothesis,
as described below. The remaining passage from arbitrary periodic GVBS
presentations to this family is separate.

\section{Support from the eventual kernel condition}
Let $h\geq0$ have positive range $R$, and suppose that
$\ker H_N(h)=G_N$ for all sufficiently large $N$.
A normalized positive state $\varphi$ satisfying
$\varphi(\tau_a(h))=0$ for every $a\in\mathbb Z$ then satisfies
$\varphi(P_{G_N})=1$ on every sufficiently large interval.
Indeed, positivity and finite dimensionality give
$H_N(h)\geq\kappa_N(1-P_{G_N})$ for some $\kappa_N>0$.
Applying $\varphi$ shows that $\varphi(1-P_{G_N})=0$.
No uniform lower bound on $\kappa_N$, translation invariance of
$\varphi$, or parent-kernel identity at the original range is required.
This implication is proved in
\path{TNLean/MPS/ParentHamiltonian/QuasiLocalGroundStateSupport.lean}.

For finitely many inequivalent primitive sectors with faithful invariant
matrices, eventual joint support also determines the entire state:
\begin{align}
  \varphi=\sum_{\alpha}t_\alpha\eta_\alpha,\qquad
  t_\alpha\geq0,\qquad\sum_{\alpha}t_\alpha=1.
  \label{eq:multiblock_supported_state_decomposition}
\end{align}
The coefficients are derived from expectations of the sector projections
on expanding centered intervals. Asymptotic orthogonality makes their total
mass tend to one; compactness supplies limiting coefficients. Compression
of each local observable identifies its expectation, and continuity extends
the identity to the completed quasi-local algebra. Translation invariance
of $\varphi$ is not assumed. This implication is proved in
\path{TNLean/MPS/ParentHamiltonian/MultiblockQuasiLocalFace.lean}.

Conversely, eventual annihilation of the joint MPS space by the
open-chain sums gives zero local expectation in every constructed
sector state. Each summand has the same expectation in such a state,
and a volume of length at least $R$ contains at least one summand.
Positivity is unnecessary for this converse implication, proved in
\path{TNLean/MPS/ParentHamiltonian/LocalParentExpectationFromOpenKernel.lean}.

The zero-energy face therefore equals the convex hull in
(\ref{eq:multiblock_supported_state_decomposition}). Since each primitive
sector state is pure among all quasi-local states, and a pure state in a
finite convex decomposition equals a constituent, the pure zero-energy
states are exactly the constructed sectors. Both equivalences are proved in
\path{TNLean/MPS/ParentHamiltonian/MultiblockQuasiLocalGroundStates.lean}.
This classification concerns the given finite primitive sector family;
more general periodic GVBS presentations are not thereby classified.

\section{One constant for finite chains and all pure ground states}
The original interaction range need not satisfy $R\geq S+1$ when the
source's eventual finite-volume kernel identity is assumed. Choose a
sufficiently large canonical range $W\geq R$ with uniform gap $\delta>0$
and with both open-chain kernels equal to the joint MPS space. Positivity
and finite dimensionality give
\begin{align}
  \kappa P_{G_W^\perp}\leq H_W(h)
  \qquad\text{for some }\kappa>0.
\end{align}
Sum this comparison over complete $W$-site windows in an $N$-site chain.
Every original range-$R$ term occurs at most $W-R+1$ times, so
\begin{align}
  \kappa H_{N,W}^{\mathrm{can}}
    \leq (W-R+1)H_N(h).
\end{align}
The coincident kernels transfer the canonical gap to the lower bound
$\kappa\delta/(W-R+1)$. A minimum with the positive gaps of the finitely
many shorter volumes gives a common constant $\gamma>0$ for every length.
The comparison and its consequence are proved in
\path{TNLean/MPS/ParentHamiltonian/Martingale/OpenInteractionRangeComparison.lean}
and \path{TNLean/MPS/ParentHamiltonian/Martingale/EventualKernelOpenGap.lean}.

Full joint projection decay and finite-range locality then give
(\ref{eq:infinite_gap_source}) in every constructed sector with this same
$\gamma$. The ground-state classification identifies every pure zero-energy
state with such a sector, so the bound holds in every pure state of the
zero-energy face. This paired finite- and infinite-volume conclusion is
proved in \path{TNLean/MPS/ParentHamiltonian/QuasiLocalGroundStateGap.lean}.
Simultaneous word span of an original, unnormalized family suffices:
normalization preserves the joint spaces at every length, and the range of
$h$ remains arbitrary and positive.

For a given inequivalent primitive family, a positive canonical interaction
with these ground states is also constructed in
\path{TNLean/MPS/ParentHamiltonian/PrimitiveGroundStateInteractionExistence.lean}.
Its range $R$ is derived, its finite-chain kernel equals the prescribed MPS
space for every $N\geq R$, and its zero-energy face and pure states have the
stated classifications. These are the existence and gap assertions for the
specified primitive representation.

\bibliographystyle{alpha}
\bibliography{references}
\end{document}
